% =====================================================================
% S4 background draft: What the Scalar Treatment Cannot Capture
% Drafted: 2026-05-29 (sub-agent 4 of 7)
% Section served: paper §4 (hinge section)
% Word target: 500--800 words
%
% Clusters drawn on:
%   C6_theory_laden -- Lavoisier/phlogiston + incommensurability subset
%                     (§1 takes the Kuhn/Hanson general-frame subset)
%   C5_echo_chambers -- Nguyen bubble-vs-chamber distinction as
%                       anticipation of the trust-topology vs
%                       trust-precision failure-mode split
%
% Cite keys used (all auto-generated from Semantic Scholar; rename
% before merging into refs.bib):
%   - kuhn1962Structure              (already in master refs.bib)
%   - Chang2012_9                    [TODO RENAME: chang2012WaterChemRev] -- ANCHOR
%   - Carrier2002_4                  [TODO RENAME: carrier2002IncommensPhlogiston]
%   - Lewowicz2011_2                 [TODO RENAME: lewowicz2011PhlogistonReferent]
%   - Mcevoy1989_7                   [TODO RENAME: mcevoy1989LavoisierLinguistic]
%   - Perrin1988_5                   [TODO RENAME: perrin1988ResearchTraditions]
%   - Boantza2008_4                  [TODO RENAME: boantza2008PhlogisticHeat]
%   - Malone1993_5                   [TODO RENAME: malone1993IncommensNoRelativism]
%   - Nguyen2018_1                   [TODO RENAME: nguyen2018BubbleChamber] -- ANCHOR
%   - Pollock2024_3                  [TODO RENAME: pollock2024ContextualDiscordance]
%   - Clinkenbeard2025_4             [TODO RENAME: clinkenbeard2025NihilisticDistrust]
%
% Anchors flagged:
%   - Lavoisier/phlogiston: Chang2012_9 (primary), Lewowicz2011_2, Carrier2002_4
%   - Nguyen bubble-vs-chamber: Nguyen2018_1 (primary, 702 citations)
%
% NOTE: Do NOT make the structural-extension argument here -- that is
% reserved for §5 (authors' move). This draft locates the scalar
% treatment's silences and points at, but does not occupy, the gap.
% =====================================================================

\section{What the Scalar Treatment Cannot Capture}
\label{sec:scalar-limits}

\paragraph{What the scalar treatment recovers.}
Before locating its limits, it is worth being explicit about what the
scalar treatment of paradigm commitment already captures. A
temperature-modulated update on a single belief variable produces an
asymmetric assimilation: the same datum is integrated by an agent who
sits near one attractor with a different effective Bayes factor than by
an agent who sits near another, so that observers in different
paradigms reach different posteriors from a shared
observation~\citep{kuhn1962Structure}. This is not a degenerate
re-description of confirmation bias; it is a principled consequence of
the cost of mind-change, and it does some of the work that Kuhn's first
face of theory-ladenness demands. To that extent, the scalar treatment
is doing real explanatory labour, and we resist the temptation to throw
it away in order to motivate what follows. The point of this section is
to mark where its explanatory reach ends.

\paragraph{Structural reorganisation: the Lavoisian case.}
The clearest pressure on a scalar treatment comes from the Chemical
Revolution. \citet{Chang2012_9} argues that the disagreement between
Lavoisier and the late phlogistonians cannot be cashed out as a
disagreement over the probability assigned to a shared proposition.
What counted as combustion, what counted as a substance, what
counted as the conservation of something, and what an experimental
yield was even evidence \emph{for}, all shifted together; the very
referents of the disputants' terms drift across the
transition~\citep{Lewowicz2011_2}. \citet{Carrier2002_4} makes the
related point that intertheoretic comparison across the phlogiston/
oxygen divide requires reconstructing each side's network of
inferential commitments rather than aligning a common variable; the
incommensurability at stake is between generative models, not between
parameter settings within one. Historical work on the linguistic and
institutional reorganisation of Lavoisian
chemistry~\citep{Mcevoy1989_7,Perrin1988_5,Boantza2008_4} reinforces
this: the revolution was a re-tooling of which questions could be
posed, not a re-weighting of answers to fixed questions. A scalar
order parameter cannot, by construction, carry this information,
because the cardinality of admissible nodes and edges is exactly what
is changing. Even philosophers sympathetic to a deflationary reading
of incommensurability concede that the residue which resists
deflation is structural rather than
quantitative~\citep{Malone1993_5}. The scalar treatment is therefore
silent on what the historians take to be the substantive content of
the revolution.

\paragraph{Failure modes the scalar treatment cannot diagnose.}
A second pressure comes from a different direction. Even if one set
aside the question of structural reorganisation and asked only after
the social pathologies a multi-agent paradigm model ought to recover,
the scalar treatment runs into a distinction it cannot draw.
\citet{Nguyen2018_1} argues that two pathologies of collective
inference -- routinely conflated in popular usage -- have entirely
different mechanisms and admit entirely different remedies. An
\emph{epistemic bubble} is a failure of \emph{informational
topology}: relevant voices are omitted from the network, coverage is
incomplete, and exposure to those voices is sufficient to dissolve
the pathology. An \emph{echo chamber} is a failure of \emph{trust
calibration}: outside voices are not absent but actively discredited,
so that exposure to them, if anything, reinforces the chamber. The
two failure modes can co-occur, but they are not the same object;
follow-up work emphasises that they decouple in
practice~\citep{Pollock2024_3} and that the trust-calibration variant
can become self-stabilising in ways that omission alone cannot
explain~\citep{Clinkenbeard2025_4}. A scalar belief variable updated
under socially-weighted observations registers both pathologies as
the same kind of thing -- a population whose marginal has collapsed
onto one attractor -- and so cannot, on its own resources,
distinguish a community that has merely failed to encounter an
alternative from one that has built the discrediting of alternatives
into its updating rule. The two patients present with the same
symptom, but they have different diseases, and the scalar model has
no diagnostic on which to separate them.

\paragraph{Locating the gap.}
The Lavoisian case and the bubble/chamber distinction press from
opposite ends on the same point. From the philosophy-of-science side,
the object being reorganised in a paradigm shift is the inferential
web, and a scalar carries none of the relevant structure. From the
social-epistemology side, the pathologies a multi-agent treatment
ought to distinguish are differentiated by their structural
locus -- topology versus trust -- and a scalar cannot mark which
locus is failing. These are not arguments against the scalar
treatment as such; they are arguments that something has to be added
to it. What that something is, and how it can be added without giving
up the dynamical traction the scalar treatment already buys, is the
business of the next section.
